Scientific claims are distributed across disciplines with different
terminologies, measurements, contexts and representational conventions. A
global synthesis can therefore fail even when every local statement is well
formed: identical terms can denote different objects, different labels can
denote compatible objects, and apparent contradictions can arise from
modality, attribution or context. This paper studies the downstream decision
of whether two already-structured scientific representations may be treated as
one statement. It does not study retrieval or raw-text extraction.

The proposed rule represents each source as a provenance-bound scientific
projection. A candidate mapping records its referent, construct, measurement,
context, polarity, modality and attribution conditions separately rather than
reducing them to a single similarity score. Compatible projections may be
integrated when the conditions required by the claim are preserved. A
load-bearing incompatibility produces a typed non-merge outcome, while
insufficient evidence leaves the relation undetermined.

We evaluate this rule on two 32-case public-reference mapping sets. The first
set motivated a prospectively frozen replication. The second set shared no case
identifiers with the first and was bound, together with its source revisions, evaluator
identities, bootstrap seed, margins and terminal rule, before confirmatory
outputs were examined. On the confirmatory holdout, coordinate-governed mapping
made no false merges, whereas flat predicate canonicalization false-merged six
of 32 cases. The paired false-merge difference was $-0.1875$ with a 95\%
bootstrap interval of $[-0.34375,-0.0625]$. The false-split difference against
an exact-coordinate conservative control was zero, satisfying the
non-inferiority guard specified in advance.

The contribution is deliberately bounded. The experiment tests mapping
decisions over already-structured projections, and all six discriminating
confirmatory cases were polarity contrasts. It
does not establish raw-text integration performance, expert construct validity
across a broader case universe, downstream answer-quality gains or the
necessity or empirical value of every coordinate. Those questions require separate prospective
evidence.

Section~\ref{sec:related} positions the decision rule relative to scientific
knowledge representation, schema matching, provenance and semantic
interoperability. Sections~\ref{sec:method}--\ref{sec:eval} define the rule and
the frozen evaluation. Section~\ref{sec:results} reports the confirmatory result
and secondary ablations. Sections~\ref{sec:limitations} and
\ref{sec:conclusion} state the boundary and conclusion.
\endinput
